\chapter{Résultats et validation}
\label{chap:results}

Une campagne complémentaire menée à l'Audio Precision après la rédaction de ce chapitre porte sur l'intermodulation, les harmoniques du secteur, le plancher de bruit large bande et la diaphonie. Faute de temps, elle n'est pas traitée ici et ses relevés sont rassemblés à l'annexe \ref{ann:imd}.

Ce chapitre confronte l'appareil terminé au cahier des charges. Le prototype est fonctionnel et assemblé dans son boîtier imprimé en 3D, pour un encombrement de 112,5 par 69,5 par 19,8 mm et une masse de 122 g. Les photographies du produit fini figurent à l'annexe~\ref{ann:photos}. Ce format tient dans une poche et satisfait la contrainte d'ergonomie posée au chapitre~\ref{chap:analyse}.

Les résultats présentés viennent d'une campagne de validation unique, complétée par des campagnes de mesure antérieures sur l'autonomie et sur le modèle d'exécution. Les valeurs commentées ici sont toutes tirées des relevés de l'annexe~\ref{ann:mesures}, qui contient les protocoles et les tableaux complets. Le chapitre ne retient que les grandeurs qui nécessitent une conclusion.

La progression suit celle du rapport. Les moyens de mesure sont d'abord présentés, puis la validation matérielle de la carte, la caractérisation de la chaîne audio, la vérification fonctionnelle du firmware et la liaison Bluetooth. Le modèle d'exécution et l'autonomie ferment la partie mesures. Le bilan final reprend chaque fonction du cahier des charges et établit ce qui a été livré.